\chapter{Circuito 1}
\label{chap:circuito_1}

Se muestra en la Figura \ref{plot:circuito_1} el circuito sobre el que se trabajará, que constituye un filtro pasa bajo
de segundo orden. Los valores de los componentes son:

\begin{itemize}
  \item $R_1: \SI{100}{\kilo\ohm}$
  \item $R_2: \SI{182}{\kilo\ohm}$
  \item $C_1: \SI{91}{\pico\farad}$
  \item $C_2: \SI{33}{\pico\farad}$
\end{itemize}

\begin{figure}[!ht]
  \centering
  \begin{tikzpicture}[transform shape]
	% Paths, nodes and wires:
	\draw (-3, -1) to[sinusoidal voltage source, l={$V_{in}$}] (-3, 1);
	\draw (-2.5, 1.75) to[american resistor, l={$R_1$}] (-0.5, 1.75);
	\draw (0, 1) to[capacitor, l_={$C_1$}] (0, -1);
	\draw (0.5, 1.75) to[american resistor, l={$R_2$}] (2.5, 1.75);
	\draw (3, 1) to[capacitor, l_={$C_2$}, v^=$V_{out}$, voltage/distance from node=0, voltage/shift=1.1, voltage=american] (3, -1);
	\draw (-3, 1) -| (-3, 1.75) -- (-2.5, 1.75);
	\draw (-0.5, 1.75) -- (0.5, 1.75);
	\draw (2.5, 1.75) -| (3, 1);
	\draw (3, -1) |- (-3, -1.75) -| (-3, -1);
	\draw (0, -1) -| (0, -1.75);
	\draw (0, 1.75) -| (0, 1);
	\node[circ] at (0, 1.75){};
	\node[circ] at (0, -1.75){};
\end{tikzpicture}

  \caption{Esquematico del circuito 1.}
  \label{plot:circuito_1}
\end{figure}

\section{Analisis}
\label{sec:analisis_circ_1}

Para el analisis se transforma el circuito a su equivalente en laplace como se ve en la Figura
\ref{plot:circuito_1_laplace} y se aplica el metodo sistematico nodal, el cual nos da las matrices de la ecuacion
\ref{eq:circuito_1_matrices}.


\begin{figure}[!ht]
  \centering
  \begin{tikzpicture}[transform shape]
	% Paths, nodes and wires:
	\draw (-3, -1) to[sinusoidal voltage source, l={$V_{in}$}] (-3, 1);
	\draw (-3, 1) -| (-3, 1.75) -- (-2.5, 1.75);
	\draw (-0.5, 1.75) -- (0.5, 1.75);
	\draw (2.5, 1.75) -| (3, 1);
	\draw (3, -1) |- (-3, -1.75) -| (-3, -1);
	\draw (0, -1) -| (0, -1.75);
	\draw (0, 1.75) -| (0, 1);
	\node[circ](N1) at (0, 1.75){} node[anchor=south] at (N1.text){$V_1$};
	\node[circ] at (0, -1.75){};
	\draw (-2.5, 1.75) to[generic, l={$\small \num{100d3}$}] (-0.5, 1.75);
	\draw (0.5, 1.75) to[generic, l={$\small \num{182d3}$}] (2.5, 1.75);
	\draw (3, 1) to[generic, l={$\frac{1}{\num{33d-12} \cdot s}$}] (3, -1);
	\draw (0, 1) to[generic, l={$\frac{1}{\num{91d-12} \cdot s}$}] (0, -1);
	\node[circ, xscale=0.1, yscale=0.1](N2) at (3, 1.75){} node[anchor=south] at (N2.text){$V_{out}$};
	\node[ground] at (0, -1.75){};
\end{tikzpicture}

  \caption{Equivalente en laplace del circuito 1.}
  \label{plot:circuito_1_laplace}
\end{figure}


\begin{equation}
  Y =
\begin{bmatrix}
  Y_{R_1} + Y_{R_2} + Y_{C_1} & -Y_{R_2} \\
  -Y_{R_2} & Y_{R_2} + Y_{C_2}
\end{bmatrix}
\qquad
I =
\begin{bmatrix}
  \frac{V_{in}(s)}{R_1} \\
0
\end{bmatrix}
  \label{eq:circuito_1_matrices}
\end{equation}

Donde:
\begin{itemize}
  \item $Y_{R_1} = \num{10d-6}$
  \item $Y_{R_2} = \num{5.49d-6}$
  \item $Y_{C_1} = \num{91d-12} \cdot s$
  \item $Y_{C_2} = \num{33d-12} \cdot s$
\end{itemize}

Al resolver la ecuación para obtener las tensiones nodales, se tiene que

\begin{equation}
V = Y^{-1} \cdot I
\end{equation}

En particular, la tensión en el nodo de salida \(V_{out}\) corresponde a la primera fila del vector \(V\), resultando:

$$
  V_{out}(s) = \frac{V_{in}(s) \cdot Y_{R_1} \cdot Y_{R_2}}{Y_{C_1} \cdot Y_{C_2} + Y_{C_1} \cdot Y_{R_1} + Y_{C_1} \cdot Y_{R_2} + Y_{R_1} \cdot Y_{R_2}}
$$

Dividiendo por $V_{in}(s)$, reemplazando los valores de las admitancias por los especificados anteriormente y
resolviendo algebraicamente, se llega finalmente la funcion de transferencia vista en la ecuacion \ref{eq:fdt-circuito-1}

\begin{align}
  H(s) &= \frac{\num{1.83d10}}{s^2 + 533134 \cdot s + \num{1.83d10}}\\
  H(s) &= \frac{\num{1.83d10}}{(s + 496214.48)\cdot(s + 36919.52)}
  \label{eq:fdt-circuito-1}
\end{align}

